\chapter{Generalizations of Legendre Symbols}

\section{Jacobi Symbols}

\section{Primes as Sums of Fourth Powers}

\begin{theorem}
  Consider an odd prime of the form $p=a^2+b^2$, where $a$ is odd and $b$ is even. Then $p$ can be written in the form
  \begin{equation*}
    p=x^4+y^4
  \end{equation*}
  if and only if $p\equiv 1\pmod{16}$ and $a\equiv 1\pmod{4}$.
\end{theorem}

\begin{proof}
  We first observe that if $p=x^4+y^4$ with $x$ odd and $y$ even, then $x^4\equiv 1\pmod{16}$ and $y^4\equiv 0\pmod{16}$ so that $p\equiv 1\pmod{16}$. Furthermore, the number $a=x^2$ is congruent to $1\pmod{4}$. This proves that the stated conditions are necessary.

  Now suppose that $p\equiv 1\pmod{16}$ and $a\equiv 1\pmod{4}$. Our goal is to show that both $a$ and $b$ are quadratic residues modulo $p$. First, we note that
  \begin{equation*}
    \legendre{a}{p}\equiv a^{\frac{p-1}{2}}\equiv (a^2)^{\frac{p-1}{4}} \equiv (-b^2)^{\frac{p-1}{4}} \equiv (b^2)^{\frac{p-1}{4}}\equiv b^{\frac{p-1}{2}}\equiv \legencre{b}{p}.
  \end{equation*}
  This shows that $a$ is a quadratic residue modulo $p$ if and only if $b$ is. To see that $a$ is a quadratic residue modulo $p$, we use the law of quadratic reciprocity for Jacobi symbols:
  \begin{equation*}
    \legendre{a}{p}\legendre{p}{a}=(-1)^{\frac{a-1}{2}\frac{p-1}{2}}.
  \end{equation*}
  Note that $\frac{p-1}{2}$ is even, so it follows that $\legendre{a}{p}=\legendre{p}{a}$. Since $p=a^2+b^2\equiv b^2\pmod{a}$, it follows that
  \begin{equation*}
    \legendre{a}{p}=\legendre{p}{a}=\legendre{b^2}{a}=\legendre{b}{a}^2=1.\qedhere
  \end{equation*}
\end{proof}
